%% =====================================================================
%%  9 Conclusion
%%  Three answers: the crossing point, the uncertainty statement, the design.
%% =====================================================================
\section{Conclusion}
\label{sec:conclusion}

Three results answer the problem posed. The first is the crossing point: from
the two measured positions at \SI{0.05}{s} and \SI{0.10}{s}, the ball reaches
the goal plane at $y=\SI{2.785}{m}$, $z=\SI{1.945}{m}$ at
$t=\SI{0.399}{s}$, \SI{0.87}{m} inside the right post and \SI{0.50}{m} below
the crossbar, with \SI{90}{\percent} intervals $y\in[2.455,3.121]$~m and
$z\in[1.617,2.276]$~m. The probability that the whole ball passes inside the
frame is $P(\text{goal}) = 0.971$, with a Monte Carlo standard error of
$0.001$; the six-parameter transverse-force model is fitted by damped
Gauss--Newton least squares to the six given coordinates at an assumed position
noise of \SI{5}{mm}, and its two transverse coefficients are effective
coefficients for whichever lateral force acts.

The second result is the uncertainty statement, and it is deliberately wider
than a statistical interval. Three ranges are reported and kept apart. The
\SI{90}{\percent} interval of the adopted model is the statistical one. The
model envelope unions the calibrated intervals of the two complete models,
$y\in[2.125,3.121]$~m and $z\in[1.617,2.276]$~m, and answers the question that
model form, not sampling, dominates the error budget. The scenario envelope
unions \SI{95}{\percent} intervals over twenty-five single-factor variants, four
joint-stress cases and one replication, $y\in[2.259,3.935]$~m and
$z\in[1.169,2.698]$~m, and is a stress range rather than a probability
statement. None of the three is a confidence interval. The calibration
experiment is the reason for preferring it. A nominal \SI{90}{\percent}
interval from the five-parameter model that the calibration tests is nearly
correctly sized in $y$ but too narrow in $z$ by a factor $\lambda_z=4.94$, a signature of the vertical force the
model does not contain, and no statistical correction applied to that interval
would repair a missing force. A reader who needs a probability statement can
have one only conditional on a model; a reader who needs a range that survives
the choice of model has the envelope, and the paper gives both while labelling
which is which.

The third result concerns the observation, not the answer. At the adopted
\SI{5}{mm} noise the two timings of the problem give a lateral uncertainty of
\SI{211}{mm}, while fixes at \SI{0.11}{s} and \SI{0.31}{s} reduce it to
\SI{12}{mm}, a factor of \num{18} (at the conservative \SI{10}{mm} the same
comparison is \SI{288}{mm} against \SI{22}{mm}). The fit and its covariance
cost \SI{8}{ms} plus \SI{0.5}{ms}, so the prediction is available about
\SI{0.11}{s} after the kick, \SI{135}{ms} before a goalkeeper can resolve
height. An independently written free-parameter analysis---a seven-parameter
spin-vector fit with its own sampler---places the crossing within \SI{3}{mm}
of the reported point and its goal probability within \num{0.002}, so the
answer is not an artefact of one parameterisation. What generalises from the
penalty kick is the ordering of the three results: the point is the cheapest
to compute, the range is the honest deliverable when the flight is only
partly observed, and the observation design controls both.
